\chapter{Purpose and Scope}
\ChapterLabel{Introduction}
\label{Section:IntroductionResponsibility}
This chapter establishes the purpose, scope, audience, and status of this specification.
\glsadd{acceptancecriterion}\glsadd{requirement}\glsadd{stakeholder}\glsadd{traceability}\glsadd{usecase}\glsadd{validation}\glsadd{verification}%


\section{Purpose}
This System Requirements and Architecture Specification records the
stakeholder needs, candidate obligations, selected architecture, and planned
verification and validation for \ProjectName. Requirements and design views
share stable references with the detailed test and evidence registers.
The authored LaTeX and editable diagram sources are authoritative; this PDF
is the integrated review publication.

\section{Scope and Boundaries}
The mission is to determine whether a measurement can be used. The system
includes the sensor and ADC0804 conversion path, clock/counter, interface
logic, FT245R USB transfer, and host-side measurement processing. The
Mission application is an external consumer. The distinction between the
whole system and the sensor processing subsystem is defined in
\hyperref[Section:MeasurementMission]{the system context}.

The current host architecture addresses macOS. The four-platform field-use
need in \UserStoryRef{US-03} remains partially covered pending
\hyperref[Question:Q-121]{Q-121}. C++ constraints remain subject to the
embedded execution allocation in \hyperref[Question:Q-116]{Q-116}.
No hardware construction, application implementation, or executed product
evidence is included in this documentation stage.

\section{Audience and Authority}
The Operator, Acquisition-Developer, Field-Engineer, and Mission-Integrator
concerns are defined in \hyperref[Table:StakeholderNeeds]{the stakeholder
register}. Michael Stoddard coordinates technical decisions and evidence
planning. Dr. Darian Muresan retains acceptance authority under
\hyperref[Table:ProjectOwners]{the owner register}. Publication layout and a
successful document build do not approve requirements or establish product
conformance.

\section{Document Status and Conventions}
This is a Draft, revision \DocumentVersion, dated \DocumentDate.
The revision identifies this document presentation; an accepted engineering
baseline remains subject to \hyperref[Question:Q-110]{Q-110}.
Template lineage 2.33 is provenance, not the project revision or a software
release version.

A statement containing \textbf{shall} expresses a candidate binding
obligation subject to its recorded lifecycle and approval conditions.
Priority uses \gls{must}, \gls{should}, \gls{could}, and \gls{would}; it is
separate from lifecycle status. TBD identifies an unresolved value or
contract. Pending verification or validation identifies missing executed
evidence, rather than a passing result.

\section{Document Organization}
\ChapterRef{SystemContext}, \ChapterRef{UserStories},
\ChapterRef{UseCases}, and \ChapterRef{Requirements} define the context,
needs, scenarios, and complete candidate requirement register.
\ChapterRef{ArchitectureOverview} and the logical, interface, interaction,
runtime, deployment, development, and implementation chapters describe the
current design and its unresolved allocations.
\ChapterRef{VerificationApproach} defines the evaluation approach;
\ChapterRef{DevelopmentPlan} owns risks, decisions, and review authority.

The appendices retain the full \hyperref[Table:ProjectTests]{test definitions},
\hyperref[Table:VerificationRegister]{verification register},
\hyperref[Table:ValidationRegister]{validation register},
\hyperref[Table:TraceabilityMatrix]{traceability matrix}, and
\ChapterRef{TestEvidence}. Each controlled statement has one source; links
connect its need, obligation, design, and planned evidence.
Repository authoring and AI guidance are maintained separately and are
excluded from this reader publication.
